\documentclass[11pt]{article}
\usepackage{preamble}
\setlength{\parskip}{4pt}

\begin{document}

\daybanner{AP Statistics \textperiodcentered\ Unit 3 \textperiodcentered\ Topic 3.13 \textperiodcentered\ TEACHER KEY}{One Result, Two Claims}

\sectionbanner{CED TETHER}

{\small
\begin{itemize}[leftmargin=1.5em,itemsep=2pt,topsep=2pt]
  \item \textbf{Skill 3.E} --- Calculate a test statistic and find a p-value, provided conditions for inference are met.
  \item \textbf{Skill 4.B} --- Interpret statistical calculations and findings to assign meaning or assess a claim.
  \item \textbf{Skill 4.E} --- Justify a claim using a decision based on significance tests.
  \item \textbf{EU VAR-6} --- The normal distribution may be used to model variation.
  \item \textbf{LO VAR-6.K} --- Calculate an appropriate test statistic for the difference of two population proportions. [Skill 3.E]
  \item \textbf{EK VAR-6.K.1} --- The test statistic for a difference in proportions is: z = [($\widehat p_1$ - $\widehat p_2$) - 0] / [$\surd$($\widehat p_c$(1-$\widehat p_c$)) $\cdot$ $\surd$(1/$n_1$ + 1/$n_2$)], where $\widehat p_c$ = ($n_1$$\widehat p_1$ + $n_2$$\widehat p_2$) / ($n_1$ + $n_2$).
  \item CLARIFYING STATEMENT: The formulas for test statistics do not appear explicitly on the AP Statistics Formula Sheet. However, these formulas do not need to be memorized, as they can be constructed based on the general test statistic formula and the standard error formulas provided on the formula sheet.
  \item \textbf{EU DAT-3} --- Significance testing allows us to make decisions about hypotheses within a particular context.
  \item \textbf{LO DAT-3.C} --- Interpret the p-value of a significance test for a difference of population proportions. [Skill 4.B]
  \item \textbf{EK DAT-3.C.1} --- An interpretation of the p-value of a significance test for a difference of two population proportions should recognize that the p-value is computed by assuming that the null hypothesis is true, i.e., by assuming that the true population proportions are equal to each other.
  \item \textbf{LO DAT-3.D} --- Justify a claim about the population based on the results of a significance test for a difference of population proportions. [Skill 4.E]
  \item \textbf{EK DAT-3.D.1} --- A formal decision explicitly compares the p-value to the significance $\alpha$. If the p-value $\leq$ $\alpha$, then reject the null hypothesis, $H_0$: $p_1$ = $p_2$, or $H_0$: $p_1$ - $p_2$ = 0. If the p-value > $\alpha$, then fail to reject the null hypothesis.
  \item \textbf{EK DAT-3.D.2} --- The results of a significance test for a difference of two population proportions can serve as the statistical reasoning to support the answer to a research question about the two populations that were sampled.
\end{itemize}
}
\textbf{Essential question:} What does a significant test against zero establish, and what practical claim remains untested?\par
\textbf{DOK-3 skill:} 4.E \textperiodcentered\ \textbf{FRQ pattern:} {\small\texttt{zero-\allowbreak{}test-\allowbreak{}vs-\allowbreak{}practical-\allowbreak{}threshold}}\par
\textbf{DOK rationale:} Part (c) coordinates the significance decision, observed effect, and a practical threshold to choose a report and explain why a test against zero does not establish the deployment requirement.\par

\textbf{Self-paced bonus sheet.} Students take it when they choose and turn it in when they finish. Everything they need is printed on the sheet. Score part (c) only, E / P / I, by hand --- never AI-graded or auto-scored. (a) and (b) are the ladder up; use them to see where a thin (c) came from.\par

\textbf{Questions to ask}
\begin{itemize}[leftmargin=1.5em,itemsep=2pt,topsep=2pt]
  \item Which null difference did this test actually use?
  \item What does 0.041 describe under that null model?
  \item How does the observed two-point gap compare with the five-point deployment standard?
  \item What additional interval or test would address the five-point claim directly?
\end{itemize}

\begin{callout}{\IconWarn\ LOOK FOR}{calloutred}
\begin{itemize}[leftmargin=1.5em,itemsep=2pt,topsep=2pt]
  \item Using unpooled standard errors in a null test of equal proportions.
  \item Interpreting the p-value as the probability that the population proportions are equal.
  \item Concluding that rejection of zero proves an improvement of at least five percentage points.
\end{itemize}
\end{callout}

\sectionbanner{SCORING PART (c) --- E / P / I}

\textbf{Expected elements}
\begin{itemize}[leftmargin=1.5em,itemsep=2pt,topsep=2pt]
  \item \textbf{judgment-1} (required) --- Uses p about 0.041 below 0.05 to reject equality.
  \item \textbf{judgment-2} (required) --- Compares the observed 2-point improvement with the 5-point target.
  \item \textbf{judgment-3} (required) --- Chooses Salma and explains that a test against zero does not establish the separate deployment threshold.
\end{itemize}

{\small\renewcommand{\arraystretch}{1.25}
\noindent\begin{tabular}{|p{0.5in}|p{5.9in}|}
\hline
\textbf{E} & Meets all three checks with correct contextual reasoning; equivalent wording is acceptable. \\ \hline
\textbf{P} & Shows at least one substantive correct check, but has an omission or error that prevents E. \\ \hline
\textbf{I} & Shows none of the three checks with substantive correct reasoning; a choice alone is insufficient. \\ \hline
\end{tabular}}

\clearpage
\focusbanner{THE PROBLEM --- annotated}

Suppose a utility has 500,000 online accounts. It randomly selects 10,000, randomly assigns 5,000 to message A and 5,000 to B, and records on-time payment. The pre-data question asks whether the population proportions differ at $\alpha=0.05$; deployment requires A to improve payment by at least 5 percentage points.\par\smallskip \textbf{Rowan:} ``I rejected equality, so A has proved an important improvement of at least 5 points. Deploy A.''\par \textbf{Salma:} ``Rejecting equality supports a nonzero difference, but does not automatically establish the separate 5-point requirement.''\par\smallskip Conditions have been verified: random selection and assignment; $10{,}000\leq50{,}000$; pooled expected counts 3,000 and 2,000 in each group.\par
\begin{center}
\textbf{On-time payment}\par\smallskip
\begin{tabular}{|l|c|c|c|}
\hline
\textbf{Text} & \textbf{n} & \textbf{Paid} & \textbf{Other} \\ \hline
\textbf{A} & 5000 & 3050 & 1950 \\ \hline
\textbf{B} & 5000 & 2950 & 2050 \\ \hline
\end{tabular}
\end{center}
\begin{firsttakebox}{FIRST TAKE --- one sentence before you work the parts}
The observed improvement is 2 percentage points and the deployment target is 5. Is that enough to deploy? Give one reason.\par
\answer{Not graded. The large samples make a two-point difference statistically detectable; students must decide what the specified deployment rule still requires.}
\end{firsttakebox}

\begin{callout}{\IconBook\ CALCULATE, INTERPRET, THEN BOUND THE CLAIM (keep on the page)}{calloutgreen}
For $H_0:p_1-p_2=0$, pool $\hat p_c=(x_1+x_2)/(n_1+n_2)$ and use
$z=\dfrac{(\hat p_1-\hat p_2)-0}{\sqrt{\hat p_c(1-\hat p_c)(1/n_1+1/n_2)}}$.
A two-sided test uses $P(|Z|\ge |z|)$. The $p$-value assumes $H_0$ is true and measures
how often random chance would produce a difference at least as extreme. If the $p$-value is at most $\alpha$, reject
$H_0$; otherwise fail to reject it. A test against zero addresses whether a difference exists;
it does not by itself establish that the difference exceeds a separate practical threshold.
Printed normal table: $P(Z\leq2.0412)=0.979387$. A right-tail area is 1 minus the left area; the two-sided area is twice that tail.
\end{callout}

\dokbadge{1}~\textbf{(a)} Calculate $\hat p_A$, $\hat p_B$, and the observed difference $\hat p_A-\hat p_B$.\par
\answer{The sample proportions are $\hat p_A=3050/5000=0.61$ and $\hat p_B=2950/5000=0.59$, so the observed difference is $0.61-0.59=0.020$, or 2 percentage points.}

\dokbadge{2}~\textbf{(b)} Using the verified conditions, calculate the pooled proportion, pooled SE, $z$, and the two-sided $p$-value for $H_0:p_A-p_B=0$.\par
\answer{Random selection and assignment support independent groups; $10{,}000\le0.10(500{,}000)=50{,}000$. The pooled proportion is $0.60$, giving expected counts 3,000 and 2,000 per group. The pooled SE is $0.00979796$, $z=2.04124$, and the two-sided $p$-value is $0.04123$. If the population proportions were equal, the chance of a sample difference at least $0.020$ in absolute value is about $0.041$. Since $0.04123<0.05$, reject $H_0$; there is evidence the proportions differ.}

\dokbadge{3}~\textbf{(c)} In 3--4 sentences, choose a report using the $p$-value, the 2-point observed gap, and the 5-point target. Explain what the test against zero establishes and why that does or does not justify deployment.\par
\answer{Salma is warranted: $0.04123<0.05$, so reject equality and report evidence of a difference. The observed improvement is 2 percentage points, below the 5-point deployment target. The test compared the difference with zero, so it does not establish an improvement of at least 5 points. Rowan's report would confuse a detectable difference with meeting the practical requirement.}

\end{document}
